% Seudocodigo completo de las etapas del pipeline, portado de algpseudocode a algorithm2e.
\subsection{Algoritmos del pipeline}
\label{sec:anexo-algoritmos}

Esta sección reúne los procedimientos completos de las etapas del pipeline que se describen en el
capítulo de Desarrollo: la percepción del archivo, la memoización de su lectura, la normalización de
cada renglón, su resolución contra el catálogo, la construcción del frente de Pareto, NSGA-II, la
selección de los kits y el frente exacto por programación dinámica.

% ---------------------------------------------------------------------
\paragraph{Percepción del archivo.}
El procedimiento recibe el archivo de una lista, lo rasteriza a 150 puntos por pulgada y lee cada
página con el modelo multimodal, validando la respuesta contra el esquema. El archivo se acepta solo
si todas sus páginas se leyeron completas y si corresponde a un único grado, contado después de
canonizar cada nombre; el tope de 40 páginas se revisa antes de gastar una llamada. Con $P$ páginas
el algoritmo hace $\Theta(P)$ llamadas, como máximo $2P$, de modo que dominan las inferencias.

\begin{algorithm}[H]
\caption{Percibir el archivo de una lista, con todas sus páginas o ninguna.}
\label{alg:anexo-percibir}
\small
\KwEntrada{$A$, archivo PDF o imagen; $\textsc{Leer}(I,j)$, lectura multimodal de la página $j$
validada contra el esquema.}
\KwSalida{$\langle ok, G, X, e\rangle$, con $ok$ verdadero solo si todas las páginas se leyeron
completas y $\lvert G\rvert \le 1$.}
\Proc{\textup{\textsc{Percibir}}$(A)$}{
  \If{$\lvert A\rvert = 0$}{
    \Return{$\langle \text{falso}, \emptyset, \langle\rangle, \textup{\textsc{ArchivoVacío}}\rangle$}\;
  }
  $\langle I_1,\dots,I_P\rangle \gets \textsc{Rasterizar}(A, 150)$\tcp*{una imagen suelta ya es una página}
  \If{$P > 40$}{
    \Return{$\langle \text{falso}, \emptyset, \langle\rangle, \textup{\textsc{DemasiadoLargo}}\rangle$}\tcp*{antes de gastar una llamada}
  }
  sean $X$ y $\Gamma$ secuencias vacías; $inc \gets \emptyset$\;
  \For{$j \gets 1$ \KwTo $P$}{
    $\langle G_j, \kappa_j, X_j\rangle \gets \textsc{Leer}(I_j, j)$\tcp*{un fallo termina con \textup{\textsc{PercepciónFallida}}}
    \If{$\kappa_j = \text{falso}$}{
      $inc \gets inc \cup \{j\}$\;
    }
    $\Gamma \gets \Gamma \cdot G_j$;\quad $X \gets X \cdot X_j$\;
  }
  $G \gets \textsc{GradosDistintos}(\Gamma)$\tcp*{con cada nombre canonizado}
  \If{$\lvert G\rvert > 1$}{
    \Return{$\langle \text{falso}, G, X, \textup{\textsc{VariosGrados}}\rangle$}\;
  }
  \If{$inc \ne \emptyset$}{
    \Return{$\langle \text{falso}, G, X, \textup{\textsc{PáginaIncompleta}}\rangle$}\;
  }
  \Return{$\langle \text{verdadero}, G, X, \emptyset\rangle$}\;
}
\end{algorithm}

\iffalse % Contenido en el repositorio: docs/algoritmos/
% ---------------------------------------------------------------------
\paragraph{Memoización de la lectura.}
Como las familias de un colegio suben el mismo archivo, la lectura aceptada se guarda en Redis
durante 30 días, con una clave que combina la huella SHA-256 del archivo y una versión del código de
lectura y del modelo. Los rechazos no se guardan, y si el almacén no responde la lectura se hace de
nuevo. Una lectura memorizada cuesta solo la huella, $\Theta(\lvert A\rvert)$.

\begin{algorithm}[H]
\caption{Memoización de la lectura de un archivo.}
\label{alg:anexo-memo}
\small
\KwEntrada{$A$, archivo; $R$, almacén clave-valor con vencimiento, posiblemente no disponible; $v$,
huella del código de lectura y del modelo.}
\Proc{\textup{\textsc{PercibirMemo}}$(A)$}{
  $k \gets v \cdot \textsc{Sha256}(A)$\;
  \If{$R$ responde \textbf{y} $R[k] \ne \bot$}{
    \Return{$R[k]$}\;
  }
  $L \gets \textsc{Percibir}(A)$\;
  \If{$L.ok$}{
    $R[k] \gets L$ con vencimiento de 30 días\tcp*{los rechazos no se guardan}
  }
  \Return{$L$}\;
}
\end{algorithm}

% ---------------------------------------------------------------------
\paragraph{Normalización de un renglón.}
La normalización recibe el texto del renglón, la cantidad de la columna si existía y el tipo y el
modo curados de la lista oficial, y devuelve una demanda con el tipo, el modo, la cantidad pedida y
el empaque que declara el texto. El tipo y el modo salen de una lista de decisión de once filas en la
que gana la primera que se cumple, y lo curado prevalece sobre lo inferido. Con $\ell$ la longitud
del renglón y un número constante de frases clave, normalizar un renglón es $\Theta(\ell)$ y una
lista, lineal en su longitud total.

\begin{algorithm}[H]
\caption{Normalizar un renglón en una demanda.}
\label{alg:anexo-normalizar}
\small
\KwEntrada{$x$, texto; $q^{\circ}$, cantidad de la columna o $\bot$; $t^{\ast}, m^{\ast}$, tipo y
modo curados o $\bot$; $D = \langle (C_r, t_r, m_r) ; 1 \le r \le 11\rangle$, lista de decisión con
$C_{11}$ siempre verdadera.}
\KwSalida{$\langle t, m, q, \upsilon\rangle$}
\Proc{\textup{\textsc{Normalizar}}$(x, q^{\circ}, t^{\ast}, m^{\ast})$}{
  $\bar x \gets \textsc{Plano}(x)$\tcp*{minúsculas, sin tildes y ``x12'' como ``x 12''}
  $r \gets \min\{\, r : \bar x \text{ contiene una frase de } C_r \,\}$\;
  $t \gets t^{\ast}$ si $t^{\ast} \ne \bot$, si no $t_r$;\quad
  $m \gets m^{\ast}$ si $m^{\ast} \ne \bot$, si no $m_r$\;
  \uIf{$t = \text{no\_comprar}$}{
    $q \gets 0$\;
  }
  \uElseIf{$q^{\circ} \ne \bot$}{
    $q \gets q^{\circ}$\tcp*{no se suma al número del texto}
  }
  \uElseIf{$\bar x$ empieza por un número $n$}{
    $q \gets n$\;
  }
  \Else{
    $q \gets 1$\;
  }
  $\upsilon \gets$ mayor $N$ de ``x $N$'', ``por $N$'' o ``caja de $N$'' sin unidad de medida
  detrás, o $\bot$\;
  \Return{$\langle t, m, q, \upsilon\rangle$}\;
}
\end{algorithm}

\fi
% ---------------------------------------------------------------------
\paragraph{Puntaje híbrido léxico-semántico.}
El procedimiento recibe el renglón ya corregido, los conceptos del catálogo con su nombre, sus
sinónimos y su centroide, y el vector del renglón si existe, y devuelve el concepto de mayor puntaje.
Si la frase del concepto empieza por la cabeza del renglón, el puntaje parte de 0,6; si la nombra en
otra posición, vale el Jaccard más el aporte del coseno, y si no la nombra solo un coseno de 0,86 o más
la rescata. El coseno se recorta en cero y solo suma, de modo que el puntaje híbrido nunca queda por
debajo del léxico. Clasificar un renglón es lineal en el catálogo, porque el coseno se calcula una vez
por concepto.

\begin{algorithm}[H]
\caption{Puntaje híbrido léxico-semántico de un renglón frente a los conceptos del catálogo.}
\label{alg:anexo-clasificar}
\small
\KwEntrada{$\tilde x$, renglón corregido; $\mathcal{C}$, conceptos con nombre, sinónimos $S(c)$ y
centroide $\bar{\mathbf{e}}(c)$; $\mathbf{e}(x)$, vector del renglón o $\bot$.}
\KwSalida{$\langle \hat c, \hat s\rangle$, concepto de mayor puntaje y su puntaje.}
\Proc{\textup{\textsc{Clasificar}}$(\tilde x, \mathcal{C}, \mathbf{e}(x))$}{
  reescribir en $\tilde x$ las series de durezas y los pares de escuadras\;
  $h \gets \textsc{Cabeza}(\tilde x)$;\quad $\langle \hat c, \hat s\rangle \gets \langle \bot, 0\rangle$\;
  \ForEach{$c \in \mathcal{C}$}{
    $\kappa \gets \max\bigl(0, \cos(\mathbf{e}(x), \bar{\mathbf{e}}(c))\bigr)$, o $0$ si falta un vector\;
    \ForEach{$y \in \{\mathrm{nom}(c)\} \cup S(c)$}{
      \uIf(\tcp*[h]{ancla, la frase empieza por lo pedido}){$h \approx y^{(1)}$}{
        $\psi \gets \min\bigl(1,\ 0{,}6 + 0{,}4\,J(\tilde x, y) + 0{,}15\,\kappa\bigr)$\;
      }
      \uElseIf{$h \approx y^{(i)}$ para algún $i$}{
        $\psi \gets \min\bigl(1,\ J(\tilde x, y) + 0{,}15\,\kappa\bigr)$\;
      }
      \uElseIf(\tcp*[h]{rescate semántico}){$\kappa \ge 0{,}86$}{
        $\psi \gets 0{,}55 + 0{,}4\,(\kappa - 0{,}86)/0{,}14$\;
      }
      \Else{
        $\psi \gets 0$\;
      }
      \If{$\psi > \hat s$}{
        $\langle \hat c, \hat s\rangle \gets \langle c, \psi\rangle$\tcp*{en empate gana el primero}
      }
    }
  }
  \Return{$\langle \hat c, \hat s\rangle$}\;
}
\end{algorithm}

% ---------------------------------------------------------------------
\paragraph{Resolución de un renglón.}
Reconocido el concepto, el procedimiento lee solo los atributos que ese concepto declara como
discriminantes y filtra sus productos. Si ninguno cumple, se abstiene en lugar de ofrecer otro, y
entre los que cumplen ofrece uno por nivel de marca, con máximo tres opciones. El estado depende de
dos umbrales: por debajo de 0,5 el renglón queda sin correspondencia, entre 0,5 y $\tau$ pasa a
revisión humana y desde $\tau$ se resuelve de forma automática. Un concepto asignado por una persona
se usa con puntaje~1, y resolver una lista es lineal en sus renglones.

\begin{algorithm}[H]
\caption{Resolver un renglón contra el catálogo y asignarle su estado.}
\label{alg:anexo-resolver}
\small
\KwEntrada{$x$, renglón; catálogo $(\mathcal{C}, \mathcal{P})$; $c^{\ast}$, concepto asignado por una
persona o $\bot$; umbral $\tau$.}
\KwSalida{$\langle e, \hat c, O\rangle$, con $e \in \{\textsc{ResueltoAuto}, \textsc{Ambiguo},
\textsc{HuecoCatálogo}, \textsc{SinMatch}\}$ y $\lvert O\rvert \le 3$.}
\Proc{\textup{\textsc{Resolver}}$(x, c^{\ast}, \tau)$}{
  $\tilde x \gets \textsc{Corregir}(x)$\tcp*{una única palabra del léxico a distancia 1}
  \uIf{$c^{\ast} \ne \bot$}{
    $\langle \hat c, \hat s\rangle \gets \langle c^{\ast}, 1\rangle$\tcp*{lo curado prevalece}
  }
  \Else{
    $\langle \hat c, \hat s\rangle \gets \textsc{Clasificar}(\tilde x, \mathcal{C}, \mathbf{e}(x))$\;
  }
  \If{$\hat c = \bot$ \textbf{o} $\hat s < 0{,}5$}{
    \Return{$\langle \textup{\textsc{SinMatch}}, \bot, \emptyset\rangle$}\;
  }
  $Imp \gets \textsc{Atributos}(\tilde x, D(\hat c))$\tcp*{solo los discriminantes del concepto}
  $K \gets \{\, p \in \mathcal{P}(\hat c) : \textsc{Cumple}(p, Imp) \,\}$\;
  \If{$K = \emptyset$}{
    \Return{$\langle \textup{\textsc{HuecoCatálogo}}, \hat c, \emptyset\rangle$}\tcp*{no se ofrece otro producto}
  }
  $O \gets$ un producto por nivel de $K$, sin personaje primero y con desempate por SKU\;
  \If{$\hat s \ge \tau$}{
    \Return{$\langle \textup{\textsc{ResueltoAuto}}, \hat c, O\rangle$}\;
  }
  \Return{$\langle \textup{\textsc{Ambiguo}}, \hat c, O\rangle$}\tcp*{pasa a revisión humana}
}
\end{algorithm}

\iffalse % Contenido en el repositorio: docs/algoritmos/
% ---------------------------------------------------------------------
\paragraph{Frente de Pareto de un hogar.}
El procedimiento recibe los candidatos admisibles de las líneas decidibles y los cuatro objetivos de
las Ecuaciones~\eqref{eq:f1f2} y~\eqref{eq:f3f4}, y construye el frente según el tamaño del espacio $\Omega$. Con
$\lvert\Omega\rvert \le 1$ no hay compromiso y decide la heurística por perfil; con
$\lvert\Omega\rvert \le 4.096$ enumera todas las asignaciones y obtiene el frente exacto, y por encima
de ese límite lo aproxima con NSGA-II sembrado. La enumeración cuesta $O(\lvert\Omega\rvert\,n)$ más
el filtro por pares, $O(\lvert U\rvert^2 M)$ con $M = 4$.

\begin{algorithm}[H]
\caption{Construir el frente de Pareto de un hogar.}
\label{alg:anexo-frente}
\small
\KwEntrada{$A_1, \dots, A_n$, candidatos admisibles de las líneas decidibles; $\mathbf{F} = (f_1, f_2,
f_3, f_4)$; semilla $\varsigma$.}
\KwSalida{Frente $\mathcal{F}$, con una asignación por vector objetivo.}
\Proc{\textup{\textsc{Frente}}$(A_1, \dots, A_n)$}{
  $\lvert\Omega\rvert \gets \prod_{i=1}^{n} \lvert A_i\rvert$\;
  \If{$n = 0$ \textbf{o} $\lvert\Omega\rvert \le 1$}{
    \Return{$\emptyset$}\tcp*{sin compromiso, decide la heurística por perfil}
  }
  \uIf{$\lvert\Omega\rvert \le 4096$}{
    $S \gets \{\, \langle \mathbf{F}(\mathbf{x}), \mathbf{x}\rangle : \mathbf{x} \in A_1 \times \cdots \times A_n \,\}$\tcp*{enumeración exacta}
  }
  \Else{
    $E \gets \textsc{Heurística}(A) \cup \textsc{Escalera}(A)$\tcp*{semillas}
    $S \gets \textsc{NSGA-II}\bigl(N(n), G(\lvert\Omega\rvert), E, \varsigma\bigr) \cup E$\;
  }
  $U \gets$ un representante de $S$ por cada vector objetivo distinto\;
  \Return{$\{\, \mathbf{x} \in U : \neg\exists\, \mathbf{y} \in U,\ \mathbf{y} \prec \mathbf{x} \,\}$
  ordenado por $(f_1, -f_2, -f_3, -f_4)$}\;
}
\end{algorithm}

% ---------------------------------------------------------------------
\paragraph{NSGA-II sembrado.}
NSGA-II mantiene $N$ soluciones, produce $N$ descendientes por selección, cruce binario simulado y
mutación polinomial, y elige la generación siguiente por rango de no dominancia y distancia de
aglomeración. La variable de cada línea es el índice entero del candidato, y la población inicial se
siembra con la heurística por perfil y con la escalera de costo marginal, de modo que el frente nunca
queda por debajo de las semillas. Su costo es $O\bigl(G(MN^2 + Nn)\bigr)$, por el ordenamiento no
dominado y la evaluación.

\begin{algorithm}[H]
\caption{NSGA-II sembrado sobre la codificación entera de índices.}
\label{alg:anexo-nsga}
\small
\KwEntrada{$N$, tamaño de la población; $G$, generaciones; $E$, semillas; $\varsigma$, semilla;
cotas $0 \le x_i \le \lvert A_i\rvert - 1$.}
\KwSalida{Aproximación del frente de Pareto.}
\Proc{\textup{\textsc{NSGA-II}}$(N, G, E, \varsigma)$}{
  $P_0 \gets E$ completada con individuos al azar generados con $\varsigma$\tcp*{si $\lvert E\rvert > N$, $N$ semillas equiespaciadas}
  evaluar $\mathbf{F}$ en $P_0$\;
  \For{$t \gets 0$ \KwTo $G-1$}{
    $Q_t \gets N$ hijos por torneo binario, cruce ($p_c = 0{,}9$, $\eta_c = 15$) y mutación
    ($\eta_m = 20$), redondeados y sin duplicados\;
    evaluar $\mathbf{F}$ en $Q_t$;\quad $R_t \gets P_t \cup Q_t$\;
    $\langle \mathcal{F}_1, \mathcal{F}_2, \dots\rangle \gets$ ordenamiento no dominado de $R_t$\tcp*{$O(M(2N)^2)$}
    $P_{t+1} \gets$ frentes completos en orden de rango hasta $N$, con el último recortado
    por distancia de aglomeración\;
  }
  \Return{los no dominados de $P_G$}\;
}
\end{algorithm}

% ---------------------------------------------------------------------
\paragraph{Selección de los kits.}
Del frente salen los kits por selección lexicográfica sobre los puntos de máxima disponibilidad. El
Ahorrador minimiza el costo, el Premium maximiza la fidelidad, el Fan maximiza el personaje, y A tu
Medida es el de mayor fidelidad cuyo costo cabe en el tope, que es el presupuesto declarado o la
mitad del camino entre el Ahorrador y el Premium. Si ningún punto trae el personaje, el Fan repite el
Ahorrador en lugar de inventarlo.

\begin{algorithm}[H]
\caption{Selección lexicográfica de los kits sobre el frente.}
\label{alg:anexo-proyectar}
\small
\KwEntrada{$\mathcal{F} \ne \emptyset$; $B$, tope declarado o $\bot$.}
\KwSalida{Asignaciones de los kits Ahorrador, Premium, Fan y A tu Medida.}
\Proc{\textup{\textsc{Proyectar}}$(\mathcal{F}, B)$}{
  $\mathcal{F}_s \gets \{\, \mathbf{x} \in \mathcal{F} : f_3(\mathbf{x}) = \max_{\mathbf{y} \in \mathcal{F}} f_3(\mathbf{y}) \,\}$\tcp*{puntos surtibles}
  $\mathbf{x}_{Ah} \gets \operatorname{lex\,min}_{\mathcal{F}_s}\,(f_1, -f_3, -f_2)$\;
  $\mathbf{x}_{Pr} \gets \operatorname{lex\,max}_{\mathcal{F}_s}\,(f_2, -f_1, f_3)$\;
  $\mathbf{x}_{Fan} \gets \operatorname{lex\,max}_{\mathcal{F}_s}\,(f_4, -f_1, f_2)$\;
  \If{$f_4(\mathbf{x}_{Fan}) = 0$}{
    $\mathbf{x}_{Fan} \gets \mathbf{x}_{Ah}$\tcp*{no se inventa un personaje}
  }
  $T \gets B$ si $B \ne \bot$, si no $\bigl\lfloor (\mathrm{total}(\mathbf{x}_{Ah}) + \mathrm{total}(\mathbf{x}_{Pr}))/2 \bigr\rfloor$\;
  $\mathbf{x}_{AM} \gets \operatorname{lex\,max}_{\{\mathbf{x} \in \mathcal{F}_s\, :\, f_1(\mathbf{x}) \le T\}}\,(f_2, -f_1, f_3)$\;
  \If{$\mathbf{x}_{AM}$ no existe \textbf{o} su total consolidado supera $T$ \textbf{o} repite
  $\mathbf{x}_{Ah}$ o $\mathbf{x}_{Pr}$}{
    $\mathbf{x}_{AM} \gets \bot$\;
  }
  \Return{$\langle \mathbf{x}_{Ah}, \mathbf{x}_{Pr}, \mathbf{x}_{Fan}, \mathbf{x}_{AM}\rangle$}\;
}
\end{algorithm}

% ---------------------------------------------------------------------
\paragraph{Frente exacto por programación dinámica.}
Como cada objetivo es una suma de términos de una sola línea y tres de ellos son contadores enteros
pequeños, el menor costo de cada terna $(a, b, p)$ de nivel de marca acumulado, líneas con precio y
líneas con personaje sale de una recurrencia línea por línea, y el frente es el conjunto de ternas no
dominadas con su costo mínimo. Tras $i$ líneas hay a lo sumo $(2i+1)(i+1)^2$ estados con $k$
transiciones cada uno, de modo que el tiempo es $O(n^4k)$ y la memoria, con el retroceso de todas las
etapas, $O(n^4)$.

\begin{algorithm}[H]
\caption{Frente exacto por programación dinámica sobre contadores enteros.}
\label{alg:anexo-pd}
\small
\KwEntrada{$A_1, \dots, A_n$; para cada candidato $c$, su nivel $\nu(c) \in \{0,1,2\}$, $\beta(c) = 1$
si tiene precio y $0$ si no, $\varphi_i(c) \in \{0,1\}$ y su costo de línea $g_i(c)$.}
\KwSalida{$C^{\ast}(a,b,p)$, menor costo de cada terna alcanzable, y el frente.}
\Proc{\textup{\textsc{FrentePD}}$(A_1, \dots, A_n)$}{
  sea $M[0..2n][0..n][0..n]$ una tabla de costos; $M \gets \langle +\infty\rangle$;\quad $M[0][0][0] \gets 0$\;
  \For{$i \gets 1$ \KwTo $n$}{
    sea $M'$ una tabla como $M$; $M' \gets \langle +\infty\rangle$\;
    \ForEach{$\langle a, b, p\rangle$ con $M[a][b][p] < +\infty$}{
      \ForEach{$c \in A_i$}{
        $\langle a', b', p'\rangle \gets \langle a + \nu(c),\ b + \beta(c),\ p + \varphi_i(c)\rangle$\;
        \If{$M[a][b][p] + g_i(c) < M'[a'][b'][p']$}{
          $M'[a'][b'][p'] \gets M[a][b][p] + g_i(c)$\tcp*{se guarda $c$ para el retroceso}
        }
      }
    }
    $M \gets M'$\;
  }
  \Return{los no dominados de $\bigl\{\, \langle M[a][b][p], \tfrac{a}{2n}, \tfrac{b}{n}, \tfrac{p}{n}\rangle : M[a][b][p] < +\infty \,\bigr\}$}\;
}
\end{algorithm}

\fi